\documentclass[11pt]{article}
\usepackage[a4paper,margin=24mm]{geometry}
\usepackage{amsmath,amssymb,amsthm,booktabs,array,hyperref,longtable,mathrsfs}
\hypersetup{colorlinks=true,linkcolor=blue,urlcolor=blue}
\newtheorem{theorem}{Theorem}
\newtheorem{lemma}[theorem]{Lemma}
\newcommand{\Z}{\mathbb Z}\newcommand{\N}{\mathbb N}
\newcommand{\Sq}{\operatorname{Sq}(\Z)}
\title{A sparse $B_4$ compiler for the ten-integer\texorpdfstring{ Con(ZFC)}{ Con(ZFC)} equation}
\author{Research supplement for ProveIt}
\date{23 October 2026}
\begin{document}\maketitle
\begin{abstract}
The pinned ten-integer construction uses base-five digit positions to
separate quartic monomials. That separation is substantially stronger than
necessary. Replacing it by an explicitly certified four-sum-distinct set
preserves existence of integer solutions through the source interface. For the same
fixed 31-variable quartic and the same 79 fixed circuit leaves, the resulting
ten-integer equation has exact total degree $69\,554\,088\,184$ and a
division-free evaluator of 1,407 binary operations, compared with degree
$589\,042\,901\,992\,797\,851\,562\,527\,064$ and 1,442 operations in the
pinned construction. This is a source-encoding improvement, not a new
ten-unknowns theorem, not a reduction in coefficient bit size, and not a
claim of global optimality. The Lucas bridge and calibrated compiler
endpoint remain explicit external dependencies.
\end{abstract}

\section{Fixed input and exact result}
The baseline is the directory \texttt{31-sun10/signed10} at ProveIt commit\\
\texttt{22e75dd6b87b1b199bb08c778d7cf075ae5aae08}.
All four native ports are
kept together, with binding \texttt{wang943-mixed-121x8801}. In particular,
the native port called \texttt{beta} must not be confused with the auxiliary
Sun scale $\beta_S$ below.

Let $r_1,\ldots,r_{22}$ be the pinned quadratic residuals in 31 positive
variables and set
\[
 P(n_1,\ldots,n_{31})=\sum_{j=1}^{22}r_j(n_1+1,\ldots,n_{31}+1)^2.
\]
This definition, its fixed ports, and all its coefficients are retained in
\texttt{sources/data}. The local verifier reconstructs the coefficient
recipes from this definition. The inherited calibration is
\begin{equation}\label{eq:endpoint}
 \mathrm{ZFC}\vdash
 \bigl[(\exists n\in\N^{31})\ P(n)=0\bigr]
 \ \longleftrightarrow\ \neg\operatorname{Con}(\mathrm{ZFC}).
\end{equation}
Equation~\eqref{eq:endpoint} is an imported mathematical endpoint, not a
claim that the Python checks establish a formal compiler-correctness proof.

Write $P=\sum_{|\alpha|\le4}p_\alpha n^\alpha$ and
\[
 c_\alpha=(4-|\alpha|)!\prod_{j=1}^{31}\alpha_j!\,p_\alpha,
 \qquad H_0=\max_\alpha|p_\alpha|.
\]
There are 402 nonzero weighted coefficient records. The data includes an
unused leading dummy coordinate $\alpha_0=0$; it is not a coded variable.
The fifth residual is 1 at the all-one assignment, so $P(0)>0$ and $H_0\ge1$.
The audited bound $|c_\alpha|+1<2^{4269855}$ permits
\[
 \beta_S=2^{4269882}>48\cdot33^4H_0.
\]
It is an even power of two and is identical to the baseline scale.

\section{Replacing exponential positions by a finite certificate}
Use the following 31 positive weights, in the order of the parent variables:
\begin{center}\small
\begin{tabular}{rrrrr}
51609&55430&68596&78030&89463\\
90233&110498&110508&128615&136498\\
211819&241519&305379&312706&314312\\
369115&417382&429683&435112&450230\\
462314&469608&490199&501539&507151\\
524944&550792&559943&573626&607799\\
610895&&&&
\end{tabular}
\end{center}
Put $t=610895$ and $L_c=4t=2443580$.

\begin{lemma}[The finite separation certificate]\label{lem:b4}
The map $\alpha\mapsto e(\alpha)=\sum_j\alpha_jw_j$ is injective on
$\{\alpha\in\N^{31}:|\alpha|\le4\}$.
\end{lemma}
\begin{proof}
Add the position zero and pad each multiset with zeros to length four.
The 52,360 nondecreasing four-tuples of indices from these 32 positions have
distinct integer sums. This is verified exhaustively by
\texttt{verify\_sun\_b4.py}, using exact integers and no random sampling.
Thus it is a finite, fully reproducible certificate of the asserted lemma.

For provenance, the candidate set comes from the Bose--Chowla construction
in the field with binary polynomial $x^{20}+x^3+1$. The generator has a full
orbit of $2^{20}-1$ nonzero elements; the subfield has 32 elements. Taking
logs of $\theta+a$, first subtracting 535448, multiplying by the unit
208418 modulo $2^{20}-1$, and then subtracting 836756 gives the displayed
positions. Equal four-term log
sums would give two monic quartics in $\theta$ with the same value; their
difference has degree at most three over the 32-element subfield, so it is
zero and the multisets agree. The construction is supplementary: the direct
finite check alone proves the particular integer claim used here.
\end{proof}

\section{A generalized sparse source lemma}
All sums and exponents in what follows have fixed integer bounds. Define
\begin{align*}
 b&=1+3f,& \mathcal B&=\beta_Sb^4,&\mathscr C&=b\mathcal B^t,\\
 \ell&=\sum_{j=1}^{31}\mathcal B^{w_j},&
 D_c&=\sum_{|\alpha|\le4}c_\alpha\mathcal B^{L_c-e(\alpha)},\\
 N_0&=\mathcal B^{t+1},&N_1&=4\mathcal B^{8t+1},&N&=N_0N_1,\\
 M_c&=N_0-1-(b-1)\ell,&
 T_c&=M_c+(\mathcal B-2)\mathcal B^{L_c}N_0,\\
 J_c&=2(1+g)^4D_c+\sum_{i=0}^{8t}\mathcal B^{i+1},&
 S_c&=g+J_cN_0,\\
 R_c&=(S_c+T_c+1)N+T_c+1,&
 X&=(N-1)R_c,&Y&=N^2.
\end{align*}
Every exponent $L_c-e(\alpha)$ is nonnegative. These are literal integer
polynomials in $f,g$, including at assignments outside the useful branch.

\begin{theorem}[Source interface]\label{thm:source}
If $P$ has a nonnegative zero, there are arbitrarily large $f>0$ with $b$ a
square power of two and $b\le g<\mathscr C$ such that
$Y\mid\binom{2X}{X}$. Conversely, for $f\ne0$, $b\in\Sq$ a power of two,
and $0\le g<2\mathscr C$, that divisibility implies a nonnegative zero of $P$.
For all $f\ne0$, $b\in\Sq$ and $0\le g<2\mathscr C$, even before assuming
$b$ is a power of two, one has
\[
 3\mid X,\qquad X>3b,\qquad Y>\max\{b,256\}.
\]
\end{theorem}
\begin{proof}
First establish the bounds for arbitrary $g$ in the indicated interval;
it is important not to assume that $g$ already has the desired digit shape.
Since $b=1+3f$ is a nonzero square and $f\ne0$, we have $b\ge4$.
Set $c=32b$. Then $2b\le c<\mathcal B$, and $1+g\le c\mathcal B^t$.
Lemma~\ref{lem:b4} and $|c_\alpha|\le24H_0$ give
\[
 0<D_c\le24H_0\sum_{j=0}^{L_c}\mathcal B^j.
\]
For the lower inequality, the leading coefficient is
$c_0=24P(0)\ge24$, whereas the absolute sum of all lower terms is less than
$\mathcal B^{L_c}$, because $24H_0<\mathcal B-1$ and all their exponents
are distinct. The scale bound also gives
$48H_0c^4\le\mathcal B-2$. Consequently
\begin{align*}
 0<J_c
 &\le(\mathcal B-2)\mathcal B^{4t}
       \sum_{j=0}^{4t}\mathcal B^j
       +\mathcal B\sum_{j=0}^{8t}\mathcal B^j\\
 &<2\mathcal B^{8t+1}<N_1.
\end{align*}
Also $0\le M_c<N_0$, $g<N_0$, and
$0\le(\mathcal B-2)\mathcal B^{L_c}<N_1$.
Thus $0\le S_c,T_c<N$. As $\beta_S\equiv b\equiv1\pmod3$, $N\equiv1\pmod3$,
giving $3\mid X$. Finally $R_c\ge N+1$, $X\ge N^2-1>3b$, and
$Y=N^2>\max\{b,256\}$.

Suppose now that $b$ is a power of two. Then $\mathcal B,N_0,N$ are powers
of two. Write $\tau_2(u,v)$ for the number of binary carries in addition.
The binary packing identity gives
\[
 Y\mid\binom{2X}{X}\quad\Longleftrightarrow\quad\tau_2(S_c,T_c)=0.
\]
For completeness, if $N=2^n$, the identity
\[
 (N-1)R_c=(S_c+T_c)N^2+(N-1-S_c)N+N-1-T_c
\]
and binary digit sums show that its digit sum is $2n-\tau_2(S_c,T_c)$.
The valuation of $\binom{2X}{X}$ is the digit sum of $X=(N-1)R_c$.
It is at least $2n$ exactly when the carry count is zero.

The low $N_0$ block and the upper block are disjoint, so the same condition
is equivalent to
\[
 \tau_2(g,M_c)=0,
 \qquad\tau_2(J_c,(\mathcal B-2)\mathcal B^{L_c})=0.
\]
The base-$\mathcal B$ digits of $M_c$ are $\mathcal B-b$ at the positions
$w_j$ and $\mathcal B-1$ elsewhere, through position $t$.
Since $g<2b\mathcal B^t\le N_0$, the first no-carry condition is exactly
\[
 g=\sum_{j=1}^{31}n_j\mathcal B^{w_j},\qquad 0\le n_j<b.
\]
Indeed a binary no-carry addition to $\mathcal B-1$ forces a zero digit,
and addition to $\mathcal B-b$ forces a digit below $b$.

For such $n$, consider
\[
 H(u)=\left(1+\sum_jn_ju^{w_j}\right)^4
       \sum_\alpha c_\alpha u^{L_c-e(\alpha)}
       =\sum_{i=0}^{8t}h_i u^i.
\]
The coefficient $h_{L_c}$ is $24P(n)$ by
Lemma~\ref{lem:b4} and the multinomial formula. All coefficients satisfy
\[
 |h_i|\le24H_0(1+\textstyle\sum_jn_j)^4<\mathcal B/2.
\]
Here distinct exponent codes ensure that, for a fixed coefficient of the
first factor, at most one coefficient of the second factor contributes to
a fixed $h_i$. The number
\[
 J'=H(\mathcal B)+(\mathcal B/2)\sum_{i=0}^{8t}\mathcal B^i
\]
therefore has genuine base-$\mathcal B$ digits
$\mathcal B/2+h_i\in[1,\mathcal B)$, and $J_c=2J'$.
The second no-carry condition is equivalent, after dividing both inputs
by two, to the no-carry addition of $J'$ and
$(\mathcal B/2-1)\mathcal B^{L_c}$. It holds exactly when
$\mathcal B/2+h_{L_c}=\mathcal B/2$, or $P(n)=0$.
This proves the reverse direction.

For the forward direction, take any zero $n$. Choose an arbitrarily large
even power of two $b$ exceeding all its coordinates; then $b\equiv1\pmod3$
and $f=(b-1)/3>0$. The encoded $g$ is less than $b\mathcal B^t$.
Since $P(0)>0$, some coordinate is nonzero, so $g\ge\mathcal B\ge b$.
The preceding equivalences establish the claimed divisibility.
\end{proof}

\section{The unchanged ten-integer bridge and finalizer}
Theorem~\ref{thm:source} supplies precisely the source interface required
by Sun's 2026 construction. Its positive bridge, reverse bridge, and full
residue interpolation are imported here, not reproved by the finite audit.
The following explicit formulas retain the pinned local three-square
quintic finalizer unchanged.

The ten unknowns, all unrestricted integers, are
$(f,g,h,k,l,w,x,y,z,m)$. All other symbols are polynomial abbreviations:
\begin{align*}
 L&=lY,&U&=2LX,&V&=4gwY,&W&=bw,\\
 A&=U(V+1),&B&=2X+1,&C&=B+(A-2)h,\\
 P_n&=U^2V,&K&=X+1+k(P_n-2),\\
 S&=2A-5,&T&=3WC-2(W^2-1),&R_2&=(P_n^2-4)K^2+4.
\end{align*}
\[
 \Delta=8\mathscr C^3gK^2-g^2\bigl(32(C-KL)^2\mathscr C^3+g^2K^2\bigr).
\]
\begin{align*}
 O&=f^2l^2\Delta,&\xi&=zx,\\
 D&=(A^2-4)C^2+4,&E&=C^2D\xi,\\
 F&=4(A^2-4)E^2+1,&G&=2(A+2)(A-2)^2E^2-1,\\
 H&=C+(B-C)F+16EFy,&I&=(G^2-1)H^2+1.
\end{align*}
Continue with
\begin{align*}
 r&=T+Sy,&\Sigma&=\xi S^5,&
 \Theta&=(r^3-rS^2-S^3)(r^2+23S^2),\\
 a_0&=b\bigl(1+(3O-4)b^2z^2\bigr),&
 a_1&=R_2,&a_2&=DFI,\\
 v_0&=1,&v_1&=a_0,&v_2&=a_0(P_n^2+K^2+3),\\
 Z&=\Sigma m-\Theta,&A_i&=\Sigma^2v_i^2a_i\quad(0\le i\le2),\\
 V_n&=Z^2+A_2-A_0-A_1,&
 W_n&=V_n^2+4Z^2A_2-4A_0A_1.
\end{align*}
The complete equation is
\begin{equation}\boxed{F_{B_4}=W_n^2-16V_n^2Z^2A_2=0.}\label{eq:final}
\end{equation}
It is division-free and has no variable exponents or unstated parameters.

\begin{theorem}[Consistency equation]
Under the inherited calibration~\eqref{eq:endpoint} and the stated Sun
bridge theorems,
\[
 \mathrm{ZFC}\vdash
 \left[\neg\exists(f,g,h,k,l,w,x,y,z,m)\in\Z^{10}\ F_{B_4}=0\right]
 \longleftrightarrow\operatorname{Con}(\mathrm{ZFC}).
\]
\end{theorem}
\begin{proof}
We record why the pinned finalizer applies to the new source rather than
silently importing its conclusion for different source polynomials.
Sun's positive construction uses only the source interface and produces
$b,R_2\in\Sq$, $S\mid T$, $C=u_B(A)$, $A>2$, and $O>0$.
Moreover $O\ge4$: $\Delta$ is a positive multiple of $g\ge b\ge4$,
and $f,l$ are nonzero.
Choose a nonzero Pell solution
$u^2-(3O-4)v^2=1$ with $b\mid v$, and set $z=v/b$.
Such divisibility follows by taking a power of a Pell unit modulo $b$;
$3O-4>0$ is nonsquare because it is $2\pmod3$.
Then $a_0$ is a square.
Full residue interpolation chooses $\xi$ divisible by this $z$ and permits
any prescribed residue of $y$ modulo $\xi$ while making $DFI$ square.
The polynomial $(q^3-q-1)(q^2+23)$ has a root modulo every nonzero integer,
so that residue can ensure $\Sigma\mid\Theta$. Put $x=\xi/z$.
All three $a_i$ are then squares, and an integer $m$ makes one factor of
the norm in~\eqref{eq:final} vanish.

Conversely, $S$ is a nonzero odd integer. Neither factor of
$(q^3-q-1)(q^2+23)$ has a rational root, so $\Theta\ne0$.
If $\Sigma=0$, the final polynomial is $\Theta^8\ne0$; hence any zero
has $\xi=zx\ne0$. Its formal norm identity is
\[
 F_{B_4}=\prod_{\varepsilon\in\{\pm1\}^3}
 \left(Z+\Sigma\sum_{i=0}^2\varepsilon_i v_i\sqrt{a_i}\right).
\]
For the actual $a_0$ and $z\ne0$, $|a_0|>1$.
Also $\sqrt{|a_1|}\le P_n^2+K^2+2$, so
\[
 |v_1|>\sqrt{|a_0|},\quad
 |v_2|>\sqrt{|a_0|}+|v_1|\sqrt{|a_1|}.
\]
A rational signed sum of these weighted roots forces all roots rational:
apply a Galois automorphism changing a root and take the largest changed
index. Its term is too large in complex absolute value to be cancelled
by earlier terms. Integer radicands with rational roots are nonnegative
integer squares. A vanishing factor therefore yields $a_i\in\Sq$ and
$\Sigma\mid\Theta$.

The coprime factors $b$ and $1+(3O-4)b^2z^2$ in $a_0$ force $b$ to be
a positive square and $O>0$. Indeed $b\equiv1\pmod3$ excludes negative
squares; a zero second factor would force $b=1$, hence $f=O=0$, and then
$1-4z^2\ne0$. If $O\le0$, the second factor is negative.
Finally, if a prime has $\nu_p(S)>\nu_p(T)$, the valuations of the two
factors of $\Theta$ are $3\nu_p(T)$ and $2\nu_p(T)$, contradicting
$S^5\mid\Theta$. Thus $S\mid T$.
These are precisely the hypotheses of Sun's reverse bridge:
$b,R_2,DFI\in\Sq$, $O>0$, $\xi\ne0$, $S\mid T$.
It returns a zero of $P$ through Theorem~\ref{thm:source}.
Negating the resulting existence equivalence and using
\eqref{eq:endpoint} finishes the proof.

The intersective quintic and separating-weight arguments above are
retained baseline lemmas. For the quintic, at an odd prime other than 23
either the cubic has a simple root, or its irreducible cubic Frobenius is
even and the discriminant $-23$ is a nonzero square; then the quadratic
has a simple root. At 23 the cubic has the simple root 3. At 2, $-23$ is
$1\pmod8$. Hensel lifting and CRT supply roots at every modulus.
\end{proof}

\section{Verification, limits, and integration}
The top degree of $A_2$ is $28464t+6766$, strictly exceeding the degrees
of $Z^2,A_0,A_1$. Thus the possible top degree of the final norm is
$4\deg A_2=113856t+27064$. Exact leading coefficients on the diagonal
specialization modulo 101, 1009, and 1000003 are respectively 95, 909,
and 668382. They are nonzero, proving attainment for the fixed tuple.

\begin{center}\begin{tabular}{lrr}
\toprule
 &Pinned baseline&This source compiler\\\midrule
Integer unknowns&10&10\\
Equations&1&1\\
Degree&$126496\cdot5^{31}+27064$&$69\,554\,088\,184$\\
Binary additions&566&484\\
Binary subtractions&25&25\\
Binary multiplications&851&898\\
Total binary operations&1442&1407\\
Fixed leaves&79&79\\\bottomrule
\end{tabular}\end{center}
All fixed powers are expanded into binary gates; reuse and constant leaves
are free, as in the baseline. The cost saving is not a claim that fewer
multiplications are used. The irregular exponentiation uses a shared
radix-16 fixed-window circuit, and the coefficient polynomial uses
homogeneous sparse Horner evaluation.

Run \texttt{python3 build\_sun\_b4.py} followed by
\texttt{python3 verify\_sun\_b4.py}. Python's standard library suffices.
The checks include exact parent coefficient reconstruction, all payload
hashes, all 52,360 four-sums, the 35-term norm identity, exact degree
certificates, and 120 modular literal-transcription tests with 5,400 named
value comparisons. The last tests are bounded regressions, not proofs of
polynomial identity. The circuit construction identities and universal
source equivalence are justified by the mathematical proof above.

The imported results are Sun 2026 Theorem 4.1 and Propositions 6.1--6.2,
whose source interface is re-established here. Their proofs use Sun 2021
Theorems 4.12--4.13. The number-theoretic proof and endpoint have not been
newly formalized in Lean or Rocq. The equation is an explicit candidate
with a complete ordinary proof relative to those stated mathematical
inputs, not an independent proof of ZFC's consistency.

\begin{thebibliography}{9}
\bibitem{sun21} Z.-W. Sun, \emph{Further results on Hilbert's Tenth Problem},
Science China Mathematics 64 (2021), 281--306.
\url{https://arxiv.org/html/1704.03504v7}.
\bibitem{sun26} Z.-W. Sun, \emph{Ten unknowns for Hilbert's tenth problem
over the integers}, preprint, version 2, 2 September 2026.
\url{https://arxiv.org/html/2609.01594v2}.
\bibitem{pin} ProveIt, immutable commit
\texttt{22e75dd6b87b1b199bb08c778d7cf075ae5aae08},
\texttt{31-sun10/signed10} under the Con(ZFC) report collection.
\url{https://github.com/VladimirReshetnikov/ProveIt/tree/22e75dd6b87b1b199bb08c778d7cf075ae5aae08}.
\end{thebibliography}
\end{document}
